\documentclass[12pt]{article}
\usepackage[utf8]{inputenc}
\usepackage[T1]{fontenc}
\usepackage[margin=1in]{geometry}
\usepackage{setspace}
\usepackage{xcolor}
\usepackage[colorlinks=true, urlcolor=blue!60!black, linkcolor=blue!60!black]{hyperref}
\setlength{\parindent}{0em}
\setlength{\parskip}{0.6em}
\renewcommand{\baselinestretch}{1.0}
\begin{document}
\begin{center}
{\LARGE \textit{United States v. E.C. Knight Co.}, 156 U.S. 1 (1895)\par}
\vspace{0.5em}
{\large Case Brief --- Two-Page Edition\par}
\end{center}

\section{Memory Jogger}
A monopoly over sugar refining was manufacturing, not commerce, so the Sherman Act could not reach it; the distinction was later repudiated (\href{https://supreme.justia.com/cases/federal/us/156/1/}{156 U.S. 1}).

\section{Detailed Case Facts}
In 1892 the American Sugar Refining Company, a New Jersey corporation already dominant in sugar refining, bought the stock of four Philadelphia refineries, including the E.C. Knight Company. Only one small Boston competitor remained, leaving the ``Sugar Trust'' with about 98\% of American refined-sugar production, a household necessity sold nationwide.

\section{Procedural History}
The United States sued in equity under the Sherman Act of 1890 to cancel the purchases as a combination to restrain and monopolize interstate trade. The Circuit Court dismissed the bill because the contracts concerned only manufacturing in Pennsylvania; the Third Circuit affirmed; the Supreme Court affirmed.

\section{Judicial Votes}
8--1. Fuller wrote for Field, Gray, Brewer, Brown, Shiras, Jackson, and White; Harlan dissented alone.

\section{Holding}
The commerce power, and thus the Sherman Act, does not reach a monopoly of manufacture within a state. Manufacturing precedes commerce and affects it only incidentally and indirectly, so it belongs to the states' police power.

\section{Analysis of Opinions}
Relying on \textit{Kidd v.\ Pearson}, Fuller defined manufacture as transformation and commerce as buying, selling, and transportation, treating state and federal spheres as exclusive. Control of manufacture controls disposition only in a secondary sense, and federal power over all goods destined for interstate trade would leave little for the states. Harlan argued that commerce includes the purchase and sale of goods meant for interstate shipment, that a 98\% trust directly burdened consumers nationwide, and that no single state could police it. \textit{Swift}, \textit{Jones \& Laughlin}, and \textit{Wickard} later abandoned the line.

\section{Examples: Future Applications}
\textbf{Same-Side:} purely local, noneconomic possession near a park cannot be aggregated (\textit{Lopez}). \textbf{Opposite-Side:} a merger giving one firm 95\% of a nationally sold chip is reachable today. \textbf{Fence-Sitter:} a ban on homemade firearms within a broad market scheme, between \textit{Lopez} and \textit{Raich}.

\section{Critique}
\textit{E.C. Knight} exemplifies formalism divorced from economic reality, as Harlan showed; its failure pushed later Courts to seek more workable limits, such as the economic/noneconomic and activity/inactivity lines.

\section{Key Quotations}
\begin{quote}
``Commerce succeeds to manufacture, and is not a part of it. \ldots That which belongs to commerce is within the jurisdiction of the United States, but that which does not belong to commerce is within the jurisdiction of the police power of the state.''
\end{quote}
\noindent --- Chief Justice Fuller, Opinion of the Court, 156 U.S. at 12.

\end{document}
